\section{What the reductions establish, and where they stop}
\label{sec:limits}

\subsection{A rejected extra elimination: trace capacity}

An attractive attempt would identify the growth-divisibility witness with
the beta-spacing witness. In a proposed shifted-horizon variant, write
\(h=t+1\), \(d=2^{19}\), and impose
\[
 F=dh(2v+1),\qquad y=F(2v+1).
\]
This would eliminate a variable, but would also force
\[
 y=\frac{F^2}{dh}.
\]
For the retained product encoding of pending labels, a state containing
two positive labels \(j_1,j_2\) already has weight
\[
 (1+j_1F)(1+j_2F)\ge(1+F)^2.
\]
Every beta modulus used at an index \(0\le i\le t\), however, would be
bounded by
\[
 1+(i+1)y\le1+hy=1+\frac{F^2}{d}<(1+F)^2.
\]
The required canonical beta residue cannot hold this state. Making the
Pell solution larger does not repair the mismatch: both sides have
quadratic dependence on \(F\), and the coefficient on the capacity side
is too small. This invalidates the inherited completeness construction
for the fused proposal. It is not, by itself, a proof that no different
trace presentation could ever satisfy the proposed equation. No
208-variable result is claimed here.

This obstruction separates two jobs that initially looked redundant.
The growth coordinate separates products of labels; the beta multiplier
must additionally accommodate the full encoded proof states. The latter
must remain independently enlargable in the construction we actually use.

\subsection{The degree-190 obstruction for this particular center}

The compact degree proof identifies a nonzero leading norm-defect term.
Consequently optimizing only the transport polynomial, while retaining
the current center and all its other definitions, cannot lower the
degree below 190 once the transport cross term is already below 190.
This is a statement about that fixed algebraic family, not a lower bound
for ten-variable Diophantine equations in general.

A possible next step is a higher-order integral center. It must be
analyzed as a whole: canceling one more leading defect term may enlarge
the center, and then the cross term with the independent carrier can
become dominant. Another direction is to change the positive quotient
chart so that the troublesome index-three case is excluded without
adding a variable or a comparably costly polynomial factor. Neither
direction has been completed in this package.

\subsection{Interpreting ``simplest''}

The principal quartic is attractive when low degree, explicit modest
coefficients, and a transparent direct logical source matter. Its
coefficient height is exactly \(2^{32}\), and all 5,486 collected
monomials are supplied. The degree-16 and degree-32 profiles trade larger
degree for fewer unrestricted integer variables. Their coefficient-height
figures are proved upper bounds, not measured exact heights.

The compact equation uses only ten positive variables, but its fixed
coefficients are very large: the largest member of its unchanged tuple
has 592,884 bits. The 119-operation count charges arithmetic gates, not
the cost of arithmetic on integers of that size. Its final collected
coefficient height and full monomial count have not been computed.
Treating these coefficients as free unit-cost literals would obscure
an important part of its description size.

No claim is made that these ideas have never appeared anywhere in the
literature. The precise novelty claim supported here is an explicit
improvement over the inspected, pinned repository constructions and their
seven incoming archives. We supply the transformations, exact equations,
and proof obligations needed to assess that claim.

Finally, a proof in ZFC of the displayed biconditionals is compatible
with the usual incompleteness phenomenon. The work proves that a zero
can be converted to a refutation and conversely; it does not prove that
either one exists, or that neither exists. No consistency assumption is
used to choose the finite coefficients.
